%% 27장 Krylov 부분공간 방법
\chapter{Krylov Subspace Methods}
\label{ch:krylov}

\abstract{행렬이 아주 크면 앞 장의 직접법도 채움 때문에 감당할 수 없게 된다. 그때는 행렬--벡터 곱만 남는다. 곱을 되풀이해 얻는 벡터 $b,Ab,A^2b,\dots$가 펼치는 공간을 Krylov 부분공간이라 한다. 이 장에서는 그 안에서 가장 좋은 근사해를 고르는 방법들을 다룬다. 대칭 양정치 행렬의 켤레 기울기법은 에너지 노름의 오차를 최소화하고 수렴 속도는 Chebyshev 다항식으로 조건수의 제곱근에 묶인다. 고유값이 모여 있으면 그 상계보다 훨씬 빨리 수렴하고, 유한 정밀도에서는 정확한 산술의 유한 종료가 무너진다는 것을 정확한 유리수 계산과 견주어 보인다. 부정 대칭 행렬의 MINRES, 비대칭 행렬의 GMRES와 BiCGSTAB, 최소제곱 문제의 LSQR과 LSMR을 차례로 유도하고, 마지막으로 반복의 수를 줄이는 전처리를 본다. stardust가 이 방법들을 Apple Accelerate에 맡기지 않고 scipy의 계산 순서대로 직접 짠 이유와, 그 결과가 스레드 수와 무관하게 같은 비트가 되게 한 방법도 함께 본다.}

\section{Krylov Subspaces}
\label{sec:krylov-subspace}

\subsection{The Polynomial Viewpoint}

$n\times n$ 정칙 행렬 $A$와 벡터 $b$에 대해
\[
\mathcal K_k(A,b)=\operatorname{span}\{b,Ab,\dots,A^{k-1}b\}
\]
를 $k$번째 \term{Krylov 부분공간}{Krylov subspace}이라 한다. 이 공간의 원소는 차수가 $k$보다 작은 다항식 $p$에 대해 $p(A)b$ 꼴로 쓸 수 있다. 그러므로 Krylov 부분공간에서 근사해를 고르는 일은 좋은 다항식을 고르는 일과 같고 이 장의 수렴 이론은 모두 다항식 근사의 문제로 바뀐다.

Krylov 부분공간은 $k$와 함께 커지다가 어느 순간 멈춘다. $A^kb$가 앞의 벡터들의 일차결합이 되는 가장 작은 $k$를 $d$라 하면 $\mathcal K_d(A,b)$는 $A$의 불변 부분공간이고 그 뒤로는 더 커지지 않는다. 이 $d$를 $b$의 $A$에 대한 \term{등급}{grade}이라 한다.

\begin{proposition}
\label{prop:krylov-contains-solution}
$A$가 정칙이고 $b$의 등급이 $d$이면 $A^{-1}b\in\mathcal K_d(A,b)$이다.
\end{proposition}

\begin{proof}
등급의 정의에 따라 $A^db=\sum_{j<d}c_jA^jb$인 계수가 있으므로, 다항식 $q(t)=t^d-\sum_{j<d}c_jt^j$에 대해 $q(A)b=0$이다. 만약 $c_0=0$이면 $q(t)=t\,\tilde q(t)$이고 $A\tilde q(A)b=0$인데, $A$가 정칙이므로 $\tilde q(A)b=0$이다. $\tilde q$의 차수가 $d-1$이므로 이것은 $b,\dots,A^{d-1}b$가 일차독립이라는 데에 어긋난다. 따라서 $c_0\ne0$이고, $q(A)b=0$을 정리하면
\[
b=\frac1{c_0}\Bigl(A^db-\sum_{j=1}^{d-1}c_jA^jb\Bigr)=A\cdot\frac1{c_0}\Bigl(A^{d-1}b-\sum_{j=1}^{d-1}c_jA^{j-1}b\Bigr)
\]
이다. 양변에 $A^{-1}$을 곱하면 $A^{-1}b$가 $b,\dots,A^{d-1}b$의 일차결합이다.
\end{proof}

그러니 정확한 산술에서는 Krylov 부분공간 방법이 많아야 $d\le n$ 걸음에 정확한 해에 닿는다. 그러나 이 성질 때문에 이 방법들이 쓰이는 것은 아니다. $n$이 수백만이면 $n$ 걸음은 의미가 없다. 실제로는 $k\ll n$에서 이미 좋은 근사가 나온다는 점이 중요하다. 이 장의 대부분은 그것이 언제 그런지를 다룬다. 행렬은 곱 $x\mapsto Ax$로만 쓰이므로 행렬을 저장할 필요조차 없고 26장의 희소 행렬--벡터 곱이나 함수로 주어진 연산자면 충분하다.

\subsection{The Arnoldi Process}

벡터 $b,Ab,A^2b,\dots$를 그대로 기저로 쓰면 안 된다. 23장의 거듭제곱법에서 보았듯이 $A^kb$는 가장 큰 고유값의 고유벡터 쪽으로 빠르게 쏠리므로 이 벡터들은 금세 거의 평행해지고 기저로서의 조건수가 지수적으로 나빠진다. 그래서 공간을 하나씩 넓힐 때마다 정규직교화한다. 이것이 \term{Arnoldi 과정}{Arnoldi process}이다\cite{arnoldi1951}. $v_1=b/\|b\|$에서 시작해 마지막 기저 벡터에 $A$를 곱하고 21장의 수정 Gram--Schmidt로 앞의 기저 벡터들의 성분을 빼서 정규화한다(\cref{alg:arnoldi}).

\begin{algorithm}[t]
\begin{algorithmic}[1]
\Require 곱 $x\mapsto Ax$, 벡터 $b\ne0$, 걸음 수 $k$
\Ensure 정규직교 열의 $V_{k+1}=[v_1,\dots,v_{k+1}]$, $(k+1)\times k$ 위 Hessenberg 행렬 $\bar H_k$
\State $v_1\gets b/\|b\|_2$
\For{$j=1,\dots,k$}
  \State $w\gets Av_j$
  \For{$i=1,\dots,j$} \Comment{수정 Gram--Schmidt}
    \State $h_{ij}\gets v_i^{\mathsf T}w$, $w\gets w-h_{ij}v_i$
  \EndFor
  \State $h_{j+1,j}\gets\|w\|_2$
  \If{$h_{j+1,j}=0$} \Return \Comment{불변 부분공간에 닿았다}
  \EndIf
  \State $v_{j+1}\gets w/h_{j+1,j}$
\EndFor
\end{algorithmic}
\caption{Arnoldi 과정}
\label{alg:arnoldi}
\end{algorithm}

알고리즘의 각 걸음은 $Av_j=\sum_{i\le j+1}h_{ij}v_i$를 뜻하므로 행렬로 모으면 \term{Arnoldi 관계}{Arnoldi relation}
\begin{equation}
\label{eq:arnoldi}
AV_k=V_{k+1}\bar H_k=V_kH_k+h_{k+1,k}v_{k+1}e_k^{\mathsf T}
\end{equation}
를 얻는다. 여기서 $H_k$는 $\bar H_k$의 위쪽 $k\times k$ 부분이고 $H_k=V_k^{\mathsf T}AV_k$이다. $H_k$는 $A$를 Krylov 부분공간에 사영한 것이다. $h_{j+1,j}=0$이면 $v_1,\dots,v_j$가 펼치는 공간이 불변이고 $j$가 등급이다. 이때 다음 절들의 방법이 정확한 해를 내므로 이것을 \term{행복한 붕괴}{lucky breakdown}라 부른다.

$k$걸음의 비용은 곱 $k$번과 내적·갱신 $O(nk^2)$이고, 기저를 모두 저장해야 하므로 메모리가 $O(nk)$이다. $k$가 커지면 이 비용이 곱보다 커진다. 그런데 $A$가 대칭이면 사정이 완전히 달라진다.

\begin{proposition}
\label{prop:lanczos-tridiagonal}
$A$가 대칭이면 $H_k$는 대칭 삼중대각 행렬이다. 따라서 Arnoldi 과정에서 $w$는 $v_{j-1}$과 $v_j$의 성분만 빼면 되고, 세 항의 점화식
\[
\beta_{j}v_{j+1}=Av_j-\alpha_jv_j-\beta_{j-1}v_{j-1},\qquad\alpha_j=v_j^{\mathsf T}Av_j
\]
를 얻는다.
\end{proposition}

\begin{proof}
$H_k=V_k^{\mathsf T}AV_k$는 $A$가 대칭이면 대칭이다. 그런데 $H_k$는 위 Hessenberg 행렬이어서 첫 부대각 아래가 $0$이고, 대칭이므로 첫 초대각 위도 $0$이다. 곧 $h_{ij}=v_i^{\mathsf T}Av_j=0$이 $i<j-1$에서 성립한다. 따라서 Arnoldi의 안쪽 반복에서 $i<j-1$인 성분은 정확한 산술에서 모두 $0$이다.
\end{proof}

이 세 항의 점화식이 \term{Lanczos 과정}{Lanczos process}이다\cite{lanczos1950}. 걸음마다 일이 $O(n)$이고 저장할 벡터가 셋뿐이다. 대칭 행렬에서는 Krylov 부분공간의 정규직교 기저가 거의 공짜로 생긴다. 다음 절의 켤레 기울기법은 이 사실을 가장 잘 살린 방법이다. 다만 유한 정밀도에서는 빼지 않은 성분이 정확히 $0$이 아니어서 직교성이 무너진다. 그 결과는 \cref{sec:cg-finite}에서 보고 고유값 계산에서 이것이 무엇을 뜻하는지는 28장에서 본다.

\section{The Conjugate Gradient Method}
\label{sec:cg}

\subsection{Minimizing the Energy Norm}

이 절에서는 $A$가 대칭 양정치라 하자. 그러면 $\|z\|_A=\sqrt{z^{\mathsf T}Az}$가 노름이고, $Ax=b$의 해 $x_\ast$는 이차 함수
\[
\phi(x)=\tfrac12x^{\mathsf T}Ax-b^{\mathsf T}x=\tfrac12\|x-x_\ast\|_A^2-\tfrac12x_\ast^{\mathsf T}Ax_\ast
\]
의 유일한 최소점이다. 둘째 등호는 전개해 보면 바로 나온다. 물리 문제에서 $\phi$는 흔히 계의 에너지이므로 $\|\cdot\|_A$를 \term{에너지 노름}{energy norm}이라 부른다. \term{켤레 기울기법}{conjugate gradient method, CG}은 $x_0=0$에서 시작해
\begin{equation}
\label{eq:cg-optimal}
x_k=\arg\min_{x\in\mathcal K_k(A,b)}\|x-x_\ast\|_A
\end{equation}
를 차례로 만드는 방법이다\cite{hestenes1952}. 초기값 $x_0$이 있으면 $b$ 대신 $r_0=b-Ax_0$에 대해 같은 일을 하면 되므로 이 절에서는 $x_0=0$으로 둔다.

\eqref{eq:cg-optimal}의 최소점은 사영으로 특징지어진다. 부분공간 위에서 $\|x-x_\ast\|_A$를 최소화하는 $x_k$는 오차 $x_\ast-x_k$가 그 부분공간과 $A$-내적에 대해 직교하는 점이다. $A(x_\ast-x_k)=b-Ax_k=r_k$이므로 이것은
\begin{equation}
\label{eq:galerkin}
r_k\perp\mathcal K_k(A,b)
\end{equation}
와 같다. 잔차가 근사 공간에 직교한다는 이 조건을 \term{Galerkin 조건}{Galerkin condition}이라 한다. 이제 이 최소점을 걸음마다 $O(n)$의 일로 갱신하는 방법을 찾아야 한다.

$p_0,\dots,p_{k-1}$이 $\mathcal K_k$의 기저이고 서로 $A$-직교, 곧 $p_i^{\mathsf T}Ap_j=0$ $(i\ne j)$이라 하자. 이런 방향들을 \term{켤레 방향}{conjugate directions}이라 한다. 그러면 $x_k=\sum_j\alpha_jp_j$로 쓸 때 $\|x-x_\ast\|_A^2$의 교차항이 사라져 각 $\alpha_j$를 따로 최소화할 수 있고, $\alpha_j=p_j^{\mathsf T}b/p_j^{\mathsf T}Ap_j$이다. 그러니 $x_{k+1}=x_k+\alpha_kp_k$이고 새 방향 하나가 더해질 때 앞의 계수는 바뀌지 않는다. 켤레 방향을 싸게 만드는 일이 남았다. 잔차 $r_k$에서 앞의 방향들의 성분을 $A$-직교화로 빼면 된다. 놀랍게도 바로 앞의 방향 하나만 빼면 충분하다.

\begin{theorem}
\label{thm:cg-short}
$r_0=b$, $p_0=r_0$으로 두고 $k=0,1,\dots$에 대해
\begin{gather*}
\alpha_k=\frac{r_k^{\mathsf T}r_k}{p_k^{\mathsf T}Ap_k},\quad x_{k+1}=x_k+\alpha_kp_k,\quad r_{k+1}=r_k-\alpha_kAp_k,\\
\beta_k=\frac{r_{k+1}^{\mathsf T}r_{k+1}}{r_k^{\mathsf T}r_k},\quad p_{k+1}=r_{k+1}+\beta_kp_k
\end{gather*}
로 정의하자. $r_k\ne0$인 동안 다음이 성립한다.
\begin{enumerate}
\item $\operatorname{span}\{r_0,\dots,r_k\}=\operatorname{span}\{p_0,\dots,p_k\}=\mathcal K_{k+1}(A,b)$.
\item 잔차는 서로 직교한다: $r_i^{\mathsf T}r_j=0$ $(i\ne j)$.
\item 방향은 서로 켤레이다: $p_i^{\mathsf T}Ap_j=0$ $(i\ne j)$.
\end{enumerate}
따라서 $x_k$는 \eqref{eq:cg-optimal}의 최소점이다.
\end{theorem}

\begin{proof}
$k$에 대한 귀납법을 쓴다. $k=0$이면 (i)은 정의이고 (ii)와 (iii)은 공허하다. $k$까지 성립한다고 하고 $k+1$을 보자.

(i): $r_{k+1}=r_k-\alpha_kAp_k$에서 $r_k\in\mathcal K_{k+1}$이고 $Ap_k\in A\mathcal K_{k+1}\subset\mathcal K_{k+2}$이므로 $r_{k+1}\in\mathcal K_{k+2}$이다. 또 $p_{k+1}=r_{k+1}+\beta_kp_k$이므로 두 span은 같다. 차원이 하나 늘었다는 것은 아래 (ii)에서 $r_{k+1}$이 앞의 잔차들과 직교하고 $0$이 아니라는 데서 나온다.

(ii): $r_{k+1}^{\mathsf T}r_j=r_k^{\mathsf T}r_j-\alpha_kp_k^{\mathsf T}Ar_j$이다. $j<k$이면 첫 항은 귀납 가정으로 $0$이고, $r_j\in\operatorname{span}\{p_0,\dots,p_j\}$이므로 둘째 항도 (iii)의 귀납 가정으로 $0$이다. $j=k$이면 $r_k=p_k-\beta_{k-1}p_{k-1}$이므로 $p_k^{\mathsf T}Ar_k=p_k^{\mathsf T}Ap_k$이고, $\alpha_k$의 정의에 의해 두 항이 상쇄된다.

(iii): $j\le k$에 대해 $p_{k+1}^{\mathsf T}Ap_j=r_{k+1}^{\mathsf T}Ap_j+\beta_kp_k^{\mathsf T}Ap_j$이다. $Ap_j=(r_j-r_{j+1})/\alpha_j$이므로 첫 항은 (ii)에 의해 $j<k$이면 $0$이고, 그러면 둘째 항도 귀납 가정으로 $0$이다. $j=k$이면 첫 항은 $-r_{k+1}^{\mathsf T}r_{k+1}/\alpha_k$이고 둘째 항은 $\beta_kp_k^{\mathsf T}Ap_k=\beta_kr_k^{\mathsf T}r_k/\alpha_k=r_{k+1}^{\mathsf T}r_{k+1}/\alpha_k$이므로 합이 $0$이다.

마지막으로 $\alpha_k$가 앞에서 구한 켤레 방향의 계수 $p_k^{\mathsf T}b/p_k^{\mathsf T}Ap_k$와 같은지 확인하자. $p_k^{\mathsf T}b=p_k^{\mathsf T}(r_k+Ax_k)=p_k^{\mathsf T}r_k$인데, $x_k$가 앞의 방향들의 일차결합이므로 $p_k^{\mathsf T}Ax_k=0$이기 때문이다. 또 $p_k^{\mathsf T}r_k=r_k^{\mathsf T}r_k+\beta_{k-1}p_{k-1}^{\mathsf T}r_k=r_k^{\mathsf T}r_k$인데, $p_{k-1}\in\mathcal K_k$이고 (i)과 (ii)에 의해 $r_k$가 $r_0,\dots,r_{k-1}$이 펼치는 $\mathcal K_k$에 직교하기 때문이다.
\end{proof}

짧은 점화식이 왜 성립하는지는 (iii)의 증명에 드러나 있다. $r_{k+1}$은 $\mathcal K_{k+1}$ 전체에 직교하는데, $A\mathcal K_k\subset\mathcal K_{k+1}$이므로 $r_{k+1}$은 $A\mathcal K_k$에도 직교하고, $A$가 대칭이므로 이것은 $r_{k+1}$이 $\mathcal K_k$와 $A$-직교한다는 말이다. 그래서 $r_{k+1}$은 이미 $p_0,\dots,p_{k-1}$과 켤레이고, 바로 앞의 $p_k$만 손보면 된다. 이 짧은 점화식은 앞 절의 Lanczos 점화식과 같은 현상이다. 실제로 CG의 잔차들을 정규화하면 Lanczos 벡터가 되고 CG는 Lanczos의 삼중대각 행렬 $T_k$에 대한 계 $T_ky=\|b\|e_1$을 $LDL^{\mathsf T}$ 분해로 한 걸음씩 푸는 것과 같다\cite{golub1989history}.

\Cref{alg:cg}는 이 점화식을 그대로 옮긴 것이다. 걸음마다 곱 하나, 내적 둘, 벡터 갱신 셋이고 저장하는 벡터는 $x,r,p,Ap$ 넷이다.

\begin{algorithm}[t]
\begin{algorithmic}[1]
\Require 대칭 양정치 $A$의 곱, $b$, 초기값 $x_0$, 허용오차 $\mathrm{rtol}$
\Ensure 근사해 $x$
\State $r\gets b-Ax_0$, $p\gets r$, $\rho\gets r^{\mathsf T}r$, $x\gets x_0$
\While{$\sqrt\rho>\mathrm{rtol}\cdot\|b\|_2$}
  \State $q\gets Ap$, $\alpha\gets\rho/(p^{\mathsf T}q)$
  \State $x\gets x+\alpha p$, $r\gets r-\alpha q$ \Comment{재귀로 갱신한 잔차}
  \State $\rho'\gets r^{\mathsf T}r$, $p\gets r+(\rho'/\rho)p$, $\rho\gets\rho'$
\EndWhile
\State \Return $x$
\end{algorithmic}
\caption{켤레 기울기법}
\label{alg:cg}
\end{algorithm}

\begin{aside}[직접법으로 태어난 반복법]
CG는 1950년 무렵 UCLA의 Hestenes와 취리히 연방공대의 Stiefel이 각자 생각해 내고 미국 국립표준국(NBS)의 수치해석연구소에서 만나 함께 발표한 방법이다\cite{hestenes1952}. 같은 연구소의 Lanczos도 같은 해에 같은 점화식에 이르렀다\cite{lanczos1952}. 당시 이 방법은 $n$걸음에 끝나는 직접법으로 소개되었다. 반올림 때문에 $n$걸음에 끝나지 않는다는 것이 곧 드러나자 Gauss 소거보다 비싼 직접법으로 여겨져 한동안 잊혔다. 1970년대 초에 큰 희소 행렬에서는 $n$보다 훨씬 적은 걸음에 충분히 좋은 근사가 나온다는 것이 다시 주목받으면서 CG는 반복법으로 부활했다. Golub과 O'Leary가 이 시기의 역사를 자세히 정리했다\cite{golub1989history}.
\end{aside}

\subsection{Convergence and Chebyshev Polynomials}

$x_k\in\mathcal K_k$이므로 $x_k=q(A)b$인 차수 $k-1$ 이하의 다항식 $q$가 있고 오차는 $x_\ast-x_k=(I-q(A)A)x_\ast=p(A)x_\ast$이다. 여기서 $p(t)=1-tq(t)$는 $p(0)=1$인 차수 $k$ 이하의 다항식이고 거꾸로 그런 다항식은 모두 이렇게 생긴다. 그러므로 \eqref{eq:cg-optimal}은
\[
\|x_\ast-x_k\|_A=\min_{p\in\mathcal P_k,\,p(0)=1}\|p(A)x_\ast\|_A
\]
와 같다. $A=Q\Lambda Q^{\mathsf T}$로 대각화하고 $x_\ast=\sum_i\xi_iq_i$로 쓰면 $\|p(A)x_\ast\|_A^2=\sum_i\lambda_ip(\lambda_i)^2\xi_i^2$이므로
\begin{equation}
\label{eq:cg-minmax}
\frac{\|x_\ast-x_k\|_A}{\|x_\ast\|_A}\le\min_{p\in\mathcal P_k,\,p(0)=1}\ \max_{\lambda\in\sigma(A)}|p(\lambda)|
\end{equation}
이다. 다시 말해 CG의 수렴은 원점에서 $1$이고 $A$의 고유값 위에서 작은 다항식이 있는지로 정해진다. 바로 다음 결과가 나온다.

\begin{corollary}
\label{cor:cg-distinct}
$A$의 서로 다른 고유값이 $m$개이면 CG는 정확한 산술에서 많아야 $m$걸음에 끝난다.
\end{corollary}

\begin{proof}
$p(t)=\prod_{i=1}^m(1-t/\lambda_i)$는 $p(0)=1$이고 모든 고유값에서 $0$이다.
\end{proof}

고유값에 대한 정보가 구간 $[\lambda_{\min},\lambda_{\max}]$뿐이라면 \eqref{eq:cg-minmax}의 최대를 구간 전체로 넓혀 최소화한다. 그 답은 Chebyshev 다항식이다. 차수 $k$의 Chebyshev 다항식은 $T_k(\cos\theta)=\cos k\theta$로 정의되며 점화식 $T_{k+1}(x)=2xT_k(x)-T_{k-1}(x)$로 보면 다항식이다. $[-1,1]$에서 $|T_k|\le1$이고, $|x|\ge1$에서는
\begin{equation}
\label{eq:chebyshev-outside}
T_k(x)=\tfrac12\Bigl[\bigl(x+\sqrt{x^2-1}\bigr)^k+\bigl(x-\sqrt{x^2-1}\bigr)^k\Bigr]
\end{equation}
이다. 이 식은 두 항이 모두 점화식을 만족하고 $k=0,1$에서 맞는다는 것으로 확인된다. Chebyshev 다항식은 29장의 보간에서 다시 중심에 놓인다.

\begin{theorem}
\label{thm:cg-chebyshev}
$A$가 대칭 양정치이고 $\kappa=\lambda_{\max}/\lambda_{\min}$이면
\[
\frac{\|x_\ast-x_k\|_A}{\|x_\ast-x_0\|_A}\le2\Bigl(\frac{\sqrt\kappa-1}{\sqrt\kappa+1}\Bigr)^k
\]
이다.
\end{theorem}

\begin{proof}
구간 $[\lambda_{\min},\lambda_{\max}]$를 $[-1,1]$로 보내는 일차 변환 $\ell(\lambda)=(\lambda_{\max}+\lambda_{\min}-2\lambda)/(\lambda_{\max}-\lambda_{\min})$을 두고
\[
p(\lambda)=\frac{T_k(\ell(\lambda))}{T_k(\ell(0))}
\]
로 놓자. $\ell(0)=(\kappa+1)/(\kappa-1)>1$이므로 분모는 양수이고 $p(0)=1$이다. 고유값 위에서는 $|T_k(\ell(\lambda))|\le1$이므로 \eqref{eq:cg-minmax}의 우변은 $1/T_k(\sigma)$ 이하이다. 여기서 $\sigma=(\kappa+1)/(\kappa-1)$이다. 그런데
\[
\sigma+\sqrt{\sigma^2-1}=\frac{\kappa+1+2\sqrt\kappa}{\kappa-1}=\frac{(\sqrt\kappa+1)^2}{(\sqrt\kappa+1)(\sqrt\kappa-1)}=\frac{\sqrt\kappa+1}{\sqrt\kappa-1}
\]
이므로 \eqref{eq:chebyshev-outside}에서 둘째 항을 버리면 $T_k(\sigma)\ge\bigl((\sqrt\kappa+1)/(\sqrt\kappa-1)\bigr)^k/2$이다. 역수를 취하면 정리의 부등식이다.
\end{proof}

이 정리를 거꾸로 읽으면 오차를 $\epsilon$배로 줄이는 데 필요한 걸음 수는 대략 $(\sqrt\kappa/2)\ln(2/\epsilon)$이다. 22장의 정규방정식이나 단순한 경사 하강법의 수렴이 $\kappa$에 비례하는 걸음을 요구하는 것과 견주면 CG는 조건수의 제곱근만 치른다. 그래서 대칭 양정치 계에서는 CG가 표준이다.

그러나 이 상계는 흔히 비관적이다. 26장의 $60\times60$ 격자 Laplace 행렬($n=3600$)은 $\kappa\approx1507.4$이어서 정리가 오차를 $10^{-8}$배로 줄이는 데 $371$걸음을 약속하지만 무작위 우변으로 실제로 돌리면 $188$걸음에 닿는다. $100$걸음에서 실제 오차 비는 $1.6\times10^{-3}$, 상계는 $1.2\times10^{-2}$이고, $150$걸음에서는 $3.7\times10^{-6}$ 대 $8.8\times10^{-4}$로 차이가 벌어진다. 걸음이 갈수록 수렴이 빨라지는 이 \term{초선형 수렴}{superlinear convergence}은 \eqref{eq:cg-minmax}의 최대가 구간 전체가 아니라 고유값 위에서만 취해지기 때문에 생긴다. CG의 다항식은 양 끝의 고유값을 근으로 잡아 먹고 나면 남은 좁은 구간의 유효 조건수로 수렴한다.

고유값이 모여 있으면 효과가 더 극적이다. $[1,2]$에 고른 $995$개의 고유값과 $10^3$부터 $10^5$까지의 고유값 $5$개를 가진 대각 행렬은 $\kappa=10^5$이어서 \cref{thm:cg-chebyshev}로는 $10^{-8}$에 $3023$걸음이 필요하다. 그러나 \cref{cor:cg-distinct}의 다항식 $5$개를 바깥 고유값에 쓰고 나머지를 $[1,2]$의 Chebyshev 다항식으로 채우면, $[1,2]$에서 $1-\lambda/\lambda_i$의 곱이 $1$ 이하이므로 $5+11=16$걸음이면 충분하다는 상계를 얻는다. 실제로 이 행렬에서 잔차를 매번 다시 직교화하며 CG를 돌리면 정확히 $16$걸음에 닿는다. 그런데 보통의 CG는 $30$걸음이 걸린다. 이 차이는 다음 절에서 다룬다.

\subsection{Conjugate Gradients in Finite Precision}
\label{sec:cg-finite}

\Cref{thm:cg-short}의 직교성은 귀납적으로 증명되었다. 반올림이 있으면 그 귀납의 각 단계에서 생긴 작은 오차가 다음 단계의 가정을 흔든다. Lanczos 과정에서 보면 새 벡터는 바로 앞의 두 벡터와만 직교화되므로 한번 생긴 옛 벡터 방향의 성분은 지워지지 않는다. Paige는 1970년대의 연구에서 이 직교성의 상실이 무작위로 일어나는 것이 아니라 어떤 Ritz 값이 고유값에 수렴하는 바로 그 순간 그 고유벡터 방향으로 일어난다는 것을 보였다\cite{paige1976error,meurant2006}. 수렴한 방향이 다시 기저에 스며들면 Lanczos는 같은 고유값을 다시 찾기 시작하고, CG는 이미 줄인 오차 성분을 다시 줄이느라 걸음을 낭비한다. 그 결과가 수렴의 \term{지연}{delay}이다.

이것을 반올림 없이 재 보려면 정확한 산술이 필요하다. Strakoš가 이 현상을 보이려고 쓴 대각 행렬\cite{strakos1991}
\[
\lambda_i=\lambda_1+\frac{i-1}{n-1}(\lambda_n-\lambda_1)\rho^{\,n-i},\qquad i=1,\dots,n
\]
에서 $n=48$, $\lambda_1=1/10$, $\lambda_n=100$, $\rho=7/8$로 두면 고유값이 모두 유리수이다. 고유값이 아래쪽에 몰리고 위쪽에 드문드문 흩어진 행렬이다. $b=(1,\dots,1)^{\mathsf T}$로 두고 CG를 Python의 유리수 산술로 정확히 돌리고 같은 계산을 binary64로 돌려 에너지 노름의 상대오차를 견주면 \cref{tab:cg-exact}와 같다.

\begin{table}[t]
\begin{sdtable}{colspec={ccc}, row{1}={font=\small\bfseries}}
\toprule
걸음 $k$ & 정확한 산술 & binary64 \\
\midrule
$10$ & $6.15\times10^{-1}$ & $6.15\times10^{-1}$ \\
$20$ & $2.08\times10^{-1}$ & $2.08\times10^{-1}$ \\
$30$ & $1.93\times10^{-2}$ & $4.79\times10^{-2}$ \\
$40$ & $5.05\times10^{-5}$ & $8.37\times10^{-3}$ \\
$47$ & $9.20\times10^{-10}$ & $1.83\times10^{-3}$ \\
$48$ & $0$ & $1.74\times10^{-3}$ \\
\bottomrule
\end{sdtable}
\caption{Strakoš 행렬($n=48$, $\kappa=1000$)에서 CG의 에너지 노름 상대오차 $\|x_\ast-x_k\|_A/\|x_\ast\|_A$. 정확한 산술은 Python의 \code{fractions}로, binary64는 numpy로 계산했다.}
\label{tab:cg-exact}
\end{table}

$20$걸음까지는 두 계산이 아홉 자리까지 같다가 $30$걸음 근처에서 갈라진다. 정확한 CG는 \cref{cor:cg-distinct}대로 $48$걸음에 오차가 정확히 $0$이 되지만, binary64의 CG는 $10^{-6}$에 $77$걸음, $10^{-10}$에 $105$걸음이 걸린다. 정확한 산술로 $44$걸음, $48$걸음이던 것이다. 같은 오차에 이르는 데 두 배 가까이, 또는 그 이상 걸린 셈이다.

그래도 binary64의 CG는 결국 수렴한다. Greenbaum은 이 사실을 정확한 명제로 만들었다. 원래 행렬의 각 고유값 근처에 고유값이 여럿 모인 더 큰 행렬을 생각하자. 유한 정밀도의 CG가 만드는 잔차 노름은 이 큰 행렬에 정확한 CG를 돌린 것과 거의 같다\cite{greenbaum1989}. 앞 절의 다항식 논증을 이 큰 행렬에 적용하면 지연은 바깥쪽의 고유값마다 근이 여러 번 필요해지는 데서 온다. 고립된 큰 고유값이 많을수록 지연이 커지는 것도 이로써 설명된다. 앞 절의 모인 고유값 예에서 $16$걸음이 $30$걸음이 된 것이 그런 경우이다. 유한 정밀도의 Lanczos와 CG에 관한 지금까지의 이해는 Meurant와 Strakoš의 개관에 정리되어 있다\cite{meurant2006}.

지연이 이보다 훨씬 심할 때도 있다. 26장의 격자 Laplace 행렬 $L$($n=3600$)의 행과 열에 $1$부터 $10^3$ 사이의 무작위 배율을 곱해 $DLD$를 만들면 $\kappa$가 $7.1\times10^7$로 커진다. 여기에 CG를 돌리면 잔차가 $10^{-8}$에 닿는 데 $44170$걸음이 걸린다. 행렬의 크기보다 열 배 넘게 많다. 이 행렬은 \cref{sec:preconditioning}에서 다시 본다.

\subsection{Stopping Criteria and the True Residual}

\Cref{alg:cg}의 $r$은 $b-Ax$를 다시 계산하지 않고 $r\gets r-\alpha q$로 갱신한 \term{재귀 잔차}{recursively updated residual}이다. 정확한 산술에서는 둘이 같지만 유한 정밀도에서는 갈라진다. 그 간격 $\|(b-Ax_k)-r_k\|$는 갱신에서 생긴 반올림의 누적이어서 대략 $u$와 $\|A\|$, 그리고 지나온 반복점 노름의 최댓값을 곱한 크기에 비례하는 데서 멈춘다\cite{greenbaum1997iterative}. 재귀 잔차는 그 뒤로도 계속 줄어들 수 있지만 참 잔차는 이 간격 아래로 내려가지 않는다. 따라서 허용오차를 이 간격보다 작게 잡으면 재귀 잔차로는 수렴했다고 판정해도 참 잔차는 판정만큼 작지 않다.

멈춤 판정의 기준도 정해야 한다. 흔한 선택은 상대 잔차 $\|r_k\|\le\mathrm{rtol}\cdot\|b\|$와 처음 잔차에 대한 상대 $\|r_k\|\le\mathrm{rtol}\cdot\|r_0\|$이다. 둘은 $x_0=0$이면 같지만 좋은 초기값을 주면 뒤의 판정은 이미 작은 잔차를 다시 $\mathrm{rtol}$배 줄이라고 요구한다. 판정의 차이는 사소해 보이지만 반복 수와 결과가 달라진다. 다른 라이브러리와 같은 답을 약속하려면 판정까지 맞추어야 한다.

\section{Symmetric Indefinite Matrices and MINRES}
\label{sec:minres}

$A$가 대칭이지만 양정치가 아니면 $\|\cdot\|_A$는 노름이 아니고 $\phi$는 아래로 유계가 아니다. CG를 그대로 돌리면 $p_k^{\mathsf T}Ap_k=0$이 되어 멈출 수 있고 그렇지 않더라도 반복점이 아무것도 최소화하지 않는다. Lanczos의 관점에서 보면 CG는 $T_ky=\|b\|e_1$을 푼다. 부정 행렬에서는 $T_k$가 특이할 수 있다.

Paige와 Saunders의 \term{MINRES}{minimal residual method}는 최소화의 대상을 잔차의 2-노름으로 바꾼다\cite{paige1975}. $x_k=V_ky$로 쓰면 Lanczos 관계 $AV_k=V_{k+1}\bar T_k$에 의해
\[
\|b-Ax_k\|_2=\|V_{k+1}(\beta_1e_1-\bar T_ky)\|_2=\|\beta_1e_1-\bar T_ky\|_2
\]
이다. 여기서 $\beta_1=\|b\|$이고 $\bar T_k$는 $(k+1)\times k$ 삼중대각 행렬이다. 마지막 등호는 $V_{k+1}$의 열이 정규직교이기 때문이다. 크기 $n$의 최소화가 크기 $k$의 최소제곱 문제로 줄었다. 이 작은 문제는 21장의 Givens 회전으로 $\bar T_k$를 위삼각으로 만들어 푸는데, $\bar T_k$는 한 걸음마다 열이 하나 늘 뿐이므로 앞에서 쓴 회전들을 새 열에 적용하고 회전 하나를 새로 만들면 된다. 회전된 우변의 마지막 성분의 절댓값 $\bar\phi_k$가 곧 잔차 노름이므로 잔차를 다시 계산하지 않고도 판정할 수 있다.

해를 갱신하는 데도 짧은 점화식이 있다. $\bar T_k$의 QR 분해 $\bar T_k=Q_k\bar R_k$에서 $R_k$는 띠 폭이 $3$인 위삼각 행렬이므로, $W_k=V_kR_k^{-1}$의 열은 $w_k=(v_k-\epsilon_kw_{k-2}-\delta_kw_{k-1})/\gamma_k$로 세 항의 점화식을 따른다. 그러면 $x_k=x_{k-1}+\phi_kw_k$이다. 결국 MINRES도 걸음마다 $O(n)$의 일과 벡터 몇 개의 저장으로 돈다.

MINRES는 잔차를 단조롭게 줄이고 부정 행렬에서도 무너지지 않지만, $\bar\phi_k$가 정확한 산술의 잔차라는 점이 약점이 된다. 유한 정밀도에서 MINRES의 해 갱신은 $W_k$를 거친다. 이 벡터들은 $R_k^{-1}$을 통해 $A$의 작은 고유값에 대한 정보를 증폭한다. Sleijpen, van der Vorst, Modersitzki는 그 결과 참 잔차와 추정 $\bar\phi_k$의 간격이 대략 $u\,\kappa(A)\,\|A\|\,\|x\|$까지 자랄 수 있음을 보였다\cite{sleijpen2000}. CG나 GMRES의 간격에 $\kappa(A)$가 하나 더 곱해진 크기이다. \Cref{tab:minres-gap}은 무작위 직교 행렬로 돌린 $100\times100$ 대칭 행렬에서 이것을 잰 것이다. 고유값은 $[1,2]$에 $90$개, $[-2,-1]$에 $8$개, 그리고 작은 고유값 $\delta$와 $-\delta/2$이고, 허용오차 $10^{-12}$로 돌렸다.

\begin{table}[t]
\begin{sdtable}{colspec={ccccc}, row{1}={font=\small\bfseries}}
\toprule
$\kappa(A)$ & 걸음 & 추정 $\bar\phi_k/\|b\|$ & 참 잔차 $\|b-Ax_k\|/\|b\|$ & GMRES의 참 잔차 \\
\midrule
$4\times10^4$ & $51$ & $5.8\times10^{-13}$ & $6.4\times10^{-11}$ & $1.2\times10^{-12}$ \\
$4\times10^6$ & $84$ & $3.3\times10^{-13}$ & $8.3\times10^{-6}$ & $4.2\times10^{-11}$ \\
$4\times10^8$ & $97$ & $2.7\times10^{-13}$ & $3.6\times10^{-2}$ & $1.7\times10^{-9}$ \\
\bottomrule
\end{sdtable}
\caption{부정 대칭 행렬에서 MINRES의 잔차 추정과 참 잔차. 마지막 열은 같은 문제를 수정 Gram--Schmidt의 완전 GMRES로 푼 것이다. 모두 numpy로 계산했다.}
\label{tab:minres-gap}
\end{table}

MINRES는 세 경우 모두 잔차가 $10^{-12}$ 아래라고 판정하고 멈추지만, $\kappa=4\times10^8$에서 참 잔차는 $3.6\times10^{-2}$이다. 같은 문제를 완전한 직교화를 하는 GMRES로 풀면 참 잔차가 판정에 훨씬 가깝다. 이 간격은 알고리즘의 결함이라기보다 짧은 점화식이 치르는 대가이다. 그래서 MINRES를 쓸 때는 끝에서 참 잔차를 다시 재는 것이 좋다.

\section{Nonsymmetric Matrices and GMRES}
\label{sec:gmres}

\subsection{The Minimal Residual Property}

$A$가 비대칭이면 Arnoldi 과정을 줄일 수 없다. Saad와 Schultz의 \term{GMRES}{generalized minimal residual method}는 MINRES의 생각을 Arnoldi 관계 \eqref{eq:arnoldi}에 그대로 옮긴다\cite{saad1986gmres}. $x_k=x_0+V_ky$에서
\[
\|b-Ax_k\|_2=\|\beta e_1-\bar H_ky\|_2,\qquad\beta=\|r_0\|_2
\]
이고, $(k+1)\times k$ Hessenberg 행렬 $\bar H_k$를 Givens 회전으로 위삼각으로 만들면 회전된 우변의 마지막 성분이 잔차 노름이다(\cref{alg:gmres}). 새 열에는 앞의 회전 $k-1$개를 적용한 뒤 부대각을 없애는 회전 하나를 만든다. 해 $y$는 판정을 통과한 뒤에 한 번만 뒤로 대입해 구한다.

\begin{algorithm}[t]
\begin{algorithmic}[1]
\Require 곱 $x\mapsto Ax$, $b$, $x_0$, 재시작 길이 $m$, 허용오차 $\mathrm{rtol}$
\Ensure 근사해 $x$
\State $x\gets x_0$
\Loop
  \State $r\gets b-Ax$, $\beta\gets\|r\|_2$
  \If{$\beta\le\mathrm{rtol}\cdot\|b\|_2$} \Return $x$
  \EndIf
  \State $v_1\gets r/\beta$, $g\gets\beta e_1$
  \For{$j=1,\dots,m$}
    \State Arnoldi의 $j$번째 걸음으로 $h_{1j},\dots,h_{j+1,j}$와 $v_{j+1}$을 만든다
    \For{$i=1,\dots,j-1$}
      \State $(h_{ij},h_{i+1,j})\gets G_i(h_{ij},h_{i+1,j})$ \Comment{앞의 회전을 새 열에}
    \EndFor
    \State $h_{jj}$와 $h_{j+1,j}$에서 $h_{j+1,j}$를 없애는 회전 $G_j$를 만든다
    \State $(g_j,g_{j+1})\gets G_j(g_j,0)$ \Comment{$|g_{j+1}|$이 잔차 노름}
    \If{$|g_{j+1}|\le\mathrm{rtol}\cdot\|b\|_2$} \textbf{break}
    \EndIf
  \EndFor
  \State 위삼각 $R y=g_{1:j}$를 풀고 $x\gets x+V_jy$
\EndLoop
\end{algorithmic}
\caption{재시작 GMRES($m$)}
\label{alg:gmres}
\end{algorithm}

Givens 회전 $G$는 $(a,b)$를 $(\sqrt{a^2+b^2},0)$으로 보내는 $c=a/r$, $s=b/r$의 회전인데, $r$을 계산할 때 8장의 넘침과 아래넘침을 피해야 한다. LAPACK의 \code{lartg}는 3.10판부터 Anderson의 안전한 척도 조정을 써서 이 계산을 가능한 한 적은 연산으로 한다\cite{anderson2017safe}.

\subsection{Convergence}

GMRES의 잔차도 다항식으로 쓸 수 있다. $r_k=p(A)r_0$이고 $p$는 $p(0)=1$인 차수 $k$ 이하의 다항식 가운데 $\|p(A)r_0\|$을 최소로 하는 것이다. 따라서 잔차는 단조롭게 줄고 정확한 산술에서는 많아야 $n$걸음에 $0$이 된다. 그러나 수렴 속도를 고유값만으로 말할 수는 없다.

\begin{theorem}
\label{thm:gmres-diag}
$A=V\Lambda V^{-1}$로 대각화되면 GMRES의 잔차는
\[
\frac{\|r_k\|_2}{\|r_0\|_2}\le\kappa_2(V)\min_{p\in\mathcal P_k,\,p(0)=1}\ \max_{\lambda\in\sigma(A)}|p(\lambda)|
\]
를 만족한다.
\end{theorem}

\begin{proof}
임의의 그런 $p$에 대해 $\|r_k\|\le\|p(A)r_0\|=\|Vp(\Lambda)V^{-1}r_0\|\le\|V\|\,\|V^{-1}\|\,\|p(\Lambda)\|\,\|r_0\|$이고, $\|p(\Lambda)\|_2=\max_i|p(\lambda_i)|$이다.
\end{proof}

$A$가 정규 행렬이면 $\kappa_2(V)=1$이어서 CG와 같은 다항식 근사의 문제가 된다. 고유값이 원점에서 떨어진 작은 원판 안에 모여 있으면 빨리 수렴하고 원점을 둘러싸면 느리다. 그러나 비정규 행렬에서는 $\kappa_2(V)$가 아주 클 수 있어 이 상계가 아무것도 말해 주지 않는다. 더 나아가 Greenbaum, Pták, Strakoš는 고유값만으로는 GMRES에 대해 아무것도 말할 수 없음을 보였다. 증가하지 않는 임의의 잔차 수열과 임의의 고유값 집합에 대해, 그 고유값을 가지면서 GMRES가 바로 그 잔차 수열을 내는 행렬과 우변이 있다\cite{greenbaum1996any}. 극단적인 예는 손으로 확인할 수 있다.

\begin{proposition}
\label{prop:gmres-shift}
$S$를 순환 이동 행렬, 곧 $Se_i=e_{i+1}$ $(i<n)$, $Se_n=e_1$이라 하고 $b=e_1$, $x_0=0$이라 하자. 그러면 GMRES의 잔차는 $k<n$에서 $\|r_k\|=1$이고 $r_n=0$이다.
\end{proposition}

\begin{proof}
$\mathcal K_k(S,e_1)=\operatorname{span}\{e_1,\dots,e_k\}$이므로 $S\mathcal K_k=\operatorname{span}\{e_2,\dots,e_{k+1}\}$이다. $k<n$이면 이 공간이 $e_1$에 직교하므로 $\|e_1-Sx\|\ge1$이고, $x=0$에서 등호가 성립한다. $k=n$이면 $S\mathcal K_n=\mathbb R^n$이다.
\end{proof}

$S$의 고유값은 $1$의 $n$제곱근으로 원점을 고르게 둘러싼다. 그러면 $p(0)=1$인 차수 $n-1$ 이하의 다항식이 고유값 위에서 작아질 수 없다. GMRES는 정확히 $n$걸음째에 한꺼번에 수렴한다. $S$는 직교 행렬이어서 조건수가 $1$인데도 그렇다.

\subsection{Restarting}

GMRES의 $k$번째 걸음은 $O(nk)$의 직교화와 $O(nk)$의 저장을 요구하므로 걸음이 많아지면 비용이 곱을 압도한다. 그래서 실제로는 $m$걸음마다 그때의 근사해에서 처음부터 다시 시작하는 \term{재시작 GMRES}{restarted GMRES}, 곧 GMRES($m$)을 쓴다(\cref{alg:gmres}). 재시작은 그동안 쌓은 Krylov 부분공간을 버리므로 최적성을 잃는다. \Cref{prop:gmres-shift}의 행렬에 $m<n$으로 재시작하면 각 주기가 $x=0$에서 아무것도 바꾸지 못하고 끝나므로 같은 주기가 영원히 되풀이된다. 이것을 \term{정체}{stagnation}라 한다.

재시작이 오히려 도움이 될 때도 있다. $50\times50$ 격자 위의 대류--확산 연산자를 상류 차분으로 이산화한 비대칭 행렬(확산 계수 $1$, 흐름 속도 $(60,30)$, $n=2500$)에 무작위 우변으로 잔차 $10^{-8}$까지 돌리면, 재시작 없는 GMRES는 $111$걸음, GMRES($40$)은 $269$걸음, GMRES($20$)은 $189$걸음, GMRES($10$)은 $154$걸음, GMRES($5$)는 $144$걸음이 걸렸다. 완전 GMRES가 걸음 수로는 가장 적지만 재시작 길이를 줄일수록 오히려 걸음 수가 줄어드는 구간이 있다. 걸음당 비용까지 생각하면 짧은 재시작이 훨씬 싸다. Embree는 재시작 길이와 수렴의 관계가 이처럼 단조롭지 않음을 작은 예로 보이고 짧은 주기는 다른 다항식들의 곱을 만들어 긴 주기가 놓치는 방향을 줄일 수 있다고 설명했다\cite{embree2003tortoise}. 좋은 재시작 길이를 미리 알 방법은 일반적으로 없다.

\section{Nonsymmetric Methods with Short Recurrences}
\label{sec:bicg}

GMRES의 비용은 직교화가 모든 기저 벡터를 거치기 때문에 자란다. 비대칭 행렬에서도 CG처럼 짧은 점화식과 최적성을 함께 가질 수는 없을까? 어떤 내적에 대한 오차를 짧은 점화식으로 최적으로 줄이는 방법이 모든 초기 벡터에 대해 존재하는지는 행렬에 달려 있다. Faber와 Manteuffel은 그런 방법이 있으려면 행렬이 그 내적에 대해 정규이고 그 수반 행렬이 행렬의 낮은 차수 다항식인 특별한 경우여야 함을 증명했다\cite{faber1984}. 그러므로 비대칭 행렬에서는 둘 가운데 하나를 포기해야 한다. GMRES는 짧은 점화식을 포기했고, 이 절의 방법들은 최적성을 포기한다.

Lanczos는 비대칭 행렬에서도 두 벡터열을 함께 만들어 세 항의 점화식을 얻는 방법을 제안했다. $\mathcal K_k(A,r_0)$의 기저 $v_j$와 $\mathcal K_k(A^{\mathsf T},\hat r_0)$의 기저 $w_j$를 서로 \term{쌍직교}{biorthogonal}, 곧 $w_i^{\mathsf T}v_j=0$ $(i\ne j)$이 되게 만들면, $W_k^{\mathsf T}AV_k$가 삼중대각이 되어 짧은 점화식이 성립한다. 이 위에 CG를 세운 것이 \term{BiCG}{biconjugate gradient}이다\cite{fletcher1976}. BiCG는 아무것도 최소화하지 않으므로 잔차가 들쭉날쭉하고, $w_k^{\mathsf T}v_k=0$인데 둘 다 $0$이 아닌 경우에는 다음 벡터를 만들 수 없다. 이것을 행복한 붕괴와 구별해 \term{심각한 붕괴}{serious breakdown}라 한다. 또 걸음마다 $A$와 $A^{\mathsf T}$를 하나씩 곱해야 한다.

Sonneveld는 BiCG의 잔차 $r_k=p_k(A)r_0$에서 $A^{\mathsf T}$ 쪽 벡터가 내적 $\hat r_0^{\mathsf T}p_k(A)^2r_0$에만 쓰인다는 것을 보고, $A^{\mathsf T}$ 없이 $p_k(A)^2r_0$을 직접 만드는 CGS를 만들었다\cite{sonneveld1989}. 수렴하면 BiCG보다 두 배 빠르지만 다항식을 제곱하므로 잔차의 요동도 제곱된다. van der Vorst의 \term{BiCGSTAB}{biconjugate gradient stabilized}은 제곱 대신 $q_k(A)p_k(A)r_0$을 만들고, $q_k(t)=(1-\omega_kt)q_{k-1}(t)$의 $\omega_k$를 걸음마다 잔차를 최소화하도록 고른다\cite{vandervorst1992}. 다시 말해 BiCG의 한 걸음 뒤에 최소 잔차의 한 걸음을 붙여 요동을 누른다. 걸음마다 $A$를 두 번 곱하고 저장하는 벡터가 몇 개뿐이어서 널리 쓰인다.

그러나 BiCGSTAB은 BiCG의 약점을 물려받는다. $1000\times1000$ 행렬의 각 행에 무작위 위치 $8$곳에 정규분포의 원소를 두고 대각에 $3$을 더한 행렬로 실험해 보자. 무작위 우변으로 잔차 $10^{-8}$까지 돌리면 BiCGSTAB은 $103$걸음(곱 $206$번)이 걸리고, 완전 GMRES는 곱 $154$번이 걸린다. 우변의 원소 하나를 한 ulp 바꾸어 스무 번 다시 돌리면 BiCGSTAB의 걸음 수는 $103$에서 $105$ 사이를 오가지만 GMRES는 언제나 $154$이다. 대각에 $3$ 대신 $2$를 더하면 고유값의 실수부가 $-0.77$에서 $4.81$로 원점을 걸친다. 이때 BiCGSTAB은 잔차가 $\|b\|$의 두 배 언저리를 맴도는 동안 $\rho_k=\hat r_0^{\mathsf T}r_k$가 $1.1\times10^3$에서 $10^{-14}$로 줄어들고, $60$걸음째에 그 몫이 NaN이 된다. 같은 행렬에서 GMRES는 곱 $614$번에 수렴한다. 반복이 길어지면 BiCGSTAB의 반올림 민감성은 더 두드러진다. 그 예는 아래의 stardust 상자에서 본다.

\section{Least Squares Problems: LSQR and LSMR}
\label{sec:lsqr}

직사각 행렬 $A\in\mathbb R^{m\times n}$의 최소제곱 문제 $\min\|b-Ax\|_2$를 22장처럼 정규방정식 $A^{\mathsf T}Ax=A^{\mathsf T}b$로 바꾸면 대칭 양반정치 계가 되므로 CG를 쓸 수 있다. $A^{\mathsf T}A$를 만들지 않고 곱 $A^{\mathsf T}(Ap)$로만 쓰는 이 방법을 CGLS라 한다. 수렴은 \cref{thm:cg-chebyshev}에서 $\kappa$를 $\kappa(A)^2$으로 바꾸어 비율이 $(\kappa(A)-1)/(\kappa(A)+1)$이 된다. 정규방정식을 만들지 않으므로 22장의 $\kappa^2$ 문제는 수렴 속도에만 나타난다.

Paige와 Saunders의 \textbf{LSQR}은 같은 반복점을 24장의 Golub--Kahan 이중대각화로 만든다\cite{paige1982lsqr,golub1965kahan}. $\beta_1u_1=b$, $\alpha_1v_1=A^{\mathsf T}u_1$에서 시작해
\[
\beta_{k+1}u_{k+1}=Av_k-\alpha_ku_k,\qquad\alpha_{k+1}v_{k+1}=A^{\mathsf T}u_{k+1}-\beta_{k+1}v_k
\]
로 두면, 정규직교 열의 $U_{k+1}$, $V_k$와 $(k+1)\times k$ 아래 이중대각 행렬 $B_k$에 대해 $AV_k=U_{k+1}B_k$이다. $V_k$의 열은 $\mathcal K_k(A^{\mathsf T}A,A^{\mathsf T}b)$를 펼친다. 그러면 $x_k=V_ky$에서 $\|b-Ax_k\|=\|\beta_1e_1-B_ky\|$이므로, MINRES와 같이 작은 최소제곱 문제를 Givens 회전으로 한 걸음씩 푼다.

\begin{proposition}
\label{prop:lsqr-cgls}
정확한 산술에서 LSQR의 $k$번째 반복점은 CGLS의 $k$번째 반복점과 같다.
\end{proposition}

\begin{proof}
두 방법 모두 반복점이 $\mathcal K_k(A^{\mathsf T}A,A^{\mathsf T}b)$에 있다. 이 공간의 $x$에 대해 $\|b-Ax\|^2=\|b-Ax_\ast\|^2+\|A(x_\ast-x)\|^2$이다. 여기서 $x_\ast$는 최소제곱 해이고, $b-Ax_\ast$가 $A$의 치역에 직교하므로 교차항이 없다. 그러므로 LSQR이 최소화하는 $\|b-Ax\|$와 CGLS가 최소화하는 $\|x_\ast-x\|_{A^{\mathsf T}A}=\|A(x_\ast-x)\|$는 같은 점에서 최소이다.
\end{proof}

두 방법은 수학적으로 같지만 실제로는 LSQR을 쓴다. Paige와 Saunders의 실험에서 유한 정밀도의 LSQR이 더 믿을 만했기 때문이다. 또 이중대각화에서 $\|A\|$와 $\kappa(A)$의 추정, 해의 노름, 정규방정식의 잔차 $\|A^{\mathsf T}r_k\|$를 거의 공짜로 얻을 수 있어 멈춤 판정이 풍부해진다. 최소제곱 문제에서는 잔차 $r_k$ 자체가 $0$으로 가지 않으므로 판정은 $\|A^{\mathsf T}r_k\|/(\|A\|\,\|r_k\|)\le\mathrm{atol}$처럼 정규방정식의 잔차로 해야 한다. 그런데 LSQR에서는 이 양이 단조롭게 줄지 않는다. Fong과 Saunders의 \textbf{LSMR}은 같은 이중대각화 위에서 $\|A^{\mathsf T}r_k\|$ 자체를 최소화한다\cite{fong2011lsmr}. LSQR이 정규방정식에 대한 CG라면 LSMR은 정규방정식에 대한 MINRES이다. 그래서 $\|A^{\mathsf T}r_k\|$가 단조롭게 줄고 일찍 멈추어도 안전하다.

두 방법 모두 Tikhonov 정칙화 $\min\|b-Ax\|^2+\lambda^2\|x\|^2$를 같은 비용으로 푼다. 이 문제는 $\begin{psmallmatrix}A\\\lambda I\end{psmallmatrix}$의 최소제곱 문제이고 그 행렬의 이중대각화는 $A$의 이중대각화에 회전 하나씩을 더해 얻을 수 있기 때문이다. 이 $\lambda$를 \term{감쇠}{damping}라 부른다.

\section{Preconditioning}
\label{sec:preconditioning}

\subsection{The Preconditioned System}

지금까지의 수렴 이론은 모두 $A$의 고유값 분포에 달려 있었다. \term{전처리}{preconditioning}는 계를 같은 해를 가지면서 고유값이 더 좋게 분포된 계로 바꾸는 일이다. $M\approx A$이고 $M^{-1}$의 곱이 싼 행렬 $M$을 골라 $M^{-1}Ax=M^{-1}b$(왼쪽 전처리)나 $AM^{-1}y=b$, $x=M^{-1}y$(오른쪽 전처리)를 푼다. 두 극단을 생각하면 쉽다. $M=I$는 공짜이지만 아무것도 바꾸지 않고, $M=A$는 한 걸음에 수렴하게 하지만 $M^{-1}$의 곱이 원래 문제이다. 좋은 전처리는 그 사이에서 $M^{-1}A$의 고유값을 모으면서 $M^{-1}$의 곱을 싸게 유지한다.

대칭 양정치 $A$에 CG를 쓰려면 전처리된 행렬도 대칭이어야 하는데, $M^{-1}A$는 일반적으로 대칭이 아니다. 그러나 $M$이 대칭 양정치이면 $M=LL^{\mathsf T}$로 쓰고 \term{나눈 전처리}{split preconditioning}
\[
\hat A\hat x=\hat b,\qquad\hat A=L^{-1}AL^{-\mathsf T},\quad\hat x=L^{\mathsf T}x,\quad\hat b=L^{-1}b
\]
에 CG를 쓸 수 있다. $\hat A$는 $M^{-1}A$와 닮은 행렬이므로 고유값이 같다. 이 계에 \cref{alg:cg}를 쓰고 변수를 원래 변수로 되돌려 쓰면 $L$이 사라지고 $M^{-1}$의 곱만 남는다.

\begin{proposition}
\label{prop:pcg}
$\hat A$에 대한 CG의 반복점 $\hat x_k$와, 다음의 \term{전처리된 켤레 기울기법}{preconditioned CG}의 반복점 $x_k$는 $\hat x_k=L^{\mathsf T}x_k$로 대응한다.
\begin{gather*}
z_k=M^{-1}r_k,\quad\alpha_k=\frac{r_k^{\mathsf T}z_k}{p_k^{\mathsf T}Ap_k},\quad x_{k+1}=x_k+\alpha_kp_k,\quad r_{k+1}=r_k-\alpha_kAp_k,\\
\beta_k=\frac{r_{k+1}^{\mathsf T}z_{k+1}}{r_k^{\mathsf T}z_k},\quad p_{k+1}=z_{k+1}+\beta_kp_k
\end{gather*}
여기서 $p_0=z_0$이다.
\end{proposition}

\begin{proof}
$\hat A$에 대한 CG의 양들을 모자로 쓰고 $\hat r_k=L^{-1}r_k$, $\hat p_k=L^{\mathsf T}p_k$로 대응시키자. 그러면 $\hat r_k^{\mathsf T}\hat r_k=r_k^{\mathsf T}L^{-\mathsf T}L^{-1}r_k=r_k^{\mathsf T}M^{-1}r_k=r_k^{\mathsf T}z_k$이고 $\hat p_k^{\mathsf T}\hat A\hat p_k=p_k^{\mathsf T}LL^{-1}AL^{-\mathsf T}L^{\mathsf T}p_k=p_k^{\mathsf T}Ap_k$이므로 $\alpha_k$와 $\beta_k$가 같다. 갱신식 $\hat x_{k+1}=\hat x_k+\alpha_k\hat p_k$에 $L^{-\mathsf T}$를, $\hat r_{k+1}=\hat r_k-\alpha_k\hat A\hat p_k$에 $L$을 곱하면 위의 $x$와 $r$의 갱신이 된다. $\hat p_{k+1}=\hat r_{k+1}+\beta_k\hat p_k$에 $L^{-\mathsf T}$를 곱하면 $L^{-\mathsf T}L^{-1}r_{k+1}=z_{k+1}$이므로 $p$의 갱신이 된다. 초기값도 같은 방식으로 대응한다.
\end{proof}

따라서 전처리된 CG의 수렴은 \cref{thm:cg-chebyshev}에서 $\kappa$를 $\kappa(M^{-1}A)$로 바꾼 것으로 묶인다. 단, 그 오차는 $\|\hat x_\ast-\hat x_k\|_{\hat A}=\|x_\ast-x_k\|_A$이므로 여전히 원래의 에너지 노름이다. 비대칭 방법에서도 같은 생각을 쓴다. 왼쪽 전처리의 GMRES는 $M^{-1}r$의 노름을 최소화하고 오른쪽 전처리는 원래 잔차의 노름을 최소화한다는 차이가 있다.

\subsection{Diagonal Preconditioning and Incomplete Factorizations}

가장 단순한 전처리는 $A$의 대각 $D$를 쓰는 \term{Jacobi 전처리}{Jacobi preconditioner}이다. 비용은 원소별 나눗셈 하나이고 행과 열의 척도가 나쁜 행렬에서 특히 효과가 크다. 앞 절의 $DLD$ 행렬이 그런 예이다. 대각으로 전처리하면 $\hat A=D_A^{-1/2}(DLD)D_A^{-1/2}$이다. 격자 Laplace 행렬의 대각은 상수 $4$이므로 $D_A=4D^2$이고 $\hat A=L/4$이다. 대각 전처리가 무작위 배율을 정확히 지우고 조건수가 $7.1\times10^7$에서 $L$의 조건수 $1507.4$로 돌아온다. CG의 걸음 수는 $44170$에서 $209$로 줄었다. 배율이 없는 $L$에 그냥 CG를 돌린 $185$걸음과 거의 같다.

대각이 이미 고른 행렬에는 Jacobi가 아무 소용이 없다. 그때는 $A$를 분해하되 채움을 버리는 \term{불완전 분해}{incomplete factorization}를 쓴다. 26장의 Cholesky를 하면서 $A$의 $0$인 자리에 생기는 채움을 모두 버리면 $A$와 같은 희소 패턴의 $L$을 얻는다. 이것이 \term{불완전 Cholesky}{incomplete Cholesky} IC(0)이다. Meijerink와 van der Vorst는 $A$가 M-행렬이면 이 분해가 무너지지 않고 전처리로 좋다는 것을 보였다\cite{meijerink1977}. $60\times60$ 격자 Laplace 행렬에서 IC(0)은 $\kappa(M^{-1}A)$를 $1507.4$에서 $134.1$로 줄이고, CG의 걸음 수를 $185$에서 $60$으로 줄인다. 비대칭 행렬에서는 같은 생각으로 불완전 LU, ILU를 쓴다. 이 밖에도 문제의 구조를 쓰는 다중격자나 영역 분할 같은 강력한 전처리가 있지만 그것은 행렬이 아니라 그 행렬을 만든 미분방정식에 관한 이야기이므로 Saad의 책\cite{saad2003iterative}에 미룬다.

\section{Iterative Methods in stardust}
\label{sec:krylov-stardust}

\begin{instardust}[\code{docs/design/sparse.md} §3.7]
Accelerate에는 CG, GMRES 등의 반복법이 있지만 stardust는 이것을 쓰지 않고 직접 짰다. 설계 문서는 멈춤 판정을 이유로 든다. Accelerate의 판정은 처음 잔차에 대한 상대이고 반복 수를 돌려주지 않으며 GMRES의 기본이 재시작 GMRES가 아닌 변형이지만, scipy는 $\|r\|\le\max(\mathrm{rtol}\cdot\|b\|,\mathrm{atol})$로 판정하고 반복 수와 상태를 알려 주며 \code{restart=20}의 재시작 GMRES를 쓴다. 한 걸음씩 부르는 인터페이스로 판정을 직접 하면 판정은 맞출 수 있지만 그 안의 직교화, 회전, 내적의 순서는 열 수 없다. 그래서 \code{cg}, \code{bicgstab}, \code{gmres}, \code{minres}, \code{lsqr}, \code{lsmr}을 scipy 1.18.1의 Python 구현에서 계산 순서까지 C++로 옮겼다. GMRES의 Givens 회전은 scipy가 부르는 LAPACK \code{lartg}의 Anderson 알고리즘을 옮긴 것이다. 전처리기 \code{M}은 scipy의 관례대로 $A^{-1}$을 어림한 것(위의 $M^{-1}$)이다.
\end{instardust}

\begin{instardust}[\code{docs/design/sparse.md} §3.7]
반복법의 비용은 곱과 내적, 벡터 갱신이다. 곱은 26장의 희소 행렬--벡터 곱이고 내적은 길이로만 정해지는 $2^{14}$개의 조각마다 부분합을 내고 조각끼리 균형 트리로 더한다. 갱신은 원소마다 \code{fma}이다. 모두 벡터의 크기만으로 일을 나누므로 같은 입력이면 스레드 수와 상관없이 반복 수와 결과의 비트가 같다(14장). 대신 numpy의 내적과 반올림 순서가 달라 끝자리가 scipy와 다를 수 있다. 설계 문서의 비교에 따르면 그 차이는 대부분 무해하지만 이 장에서 본 반올림 민감성이 그대로 드러나는 문제가 있다. 무작위 비대칭 행렬의 BiCGSTAB에서는 우변을 상대적으로 $10^{-16}$만큼 바꾼 scipy끼리도 반복 수가 $54$에서 $71$까지 퍼졌다. 그래서 그런 문제에서는 반복 수가 scipy 자신의 퍼짐 안에 있는지를 본다. 또 조건수가 $10^{10}$인 대칭 행렬에서 \code{minres}는 허용오차 $10^{-10}$로 수렴했다고 보고하지만 참 상대 잔차는 $7\times10^{-2}$였고 scipy도 같았다. \cref{tab:minres-gap}의 현상이다. 설계 문서는 이것을 결함으로 고치지 않고 정확도 계약에 적어 두었다.
\end{instardust}

\section{Concluding Remarks}

Krylov 부분공간 방법은 행렬--벡터 곱 하나로 할 수 있는 최선을 다항식 근사 문제로 바꾸어 다룬다. 대칭 양정치 행렬에서는 Lanczos의 세 항 점화식 덕분에 짧은 점화식과 최적성을 함께 얻는 CG가 있다. 그 수렴은 Chebyshev 다항식이 말해 주는 $\sqrt\kappa$보다 대개 빠르다. 대칭성을 잃으면 GMRES처럼 최적성을 지키고 비용을 치르거나 BiCGSTAB처럼 비용을 아끼고 최적성과 안정성을 내주어야 한다. 어느 쪽이든 유한 정밀도는 정확한 산술에서 얻은 결론을 바꾸어 놓는다. 직교성이 무너지면 수렴이 늦어지고 짧은 점화식 때문에 참 잔차와 추정 잔차가 갈라진다. 다음 장에서는 같은 Krylov 부분공간과 Lanczos, Arnoldi 과정을 고유값 문제에 쓴다. 거기서는 직교성의 상실이 가짜 고유값이라는 더 눈에 띄는 모습으로 나타난다.

Krylov 부분공간 방법의 교과서로는 Saad\cite{saad2003iterative}와 Greenbaum\cite{greenbaum1997iterative}의 책이 표준이다. 수렴 이론과 유한 정밀도의 행동을 깊이 다룬 것으로는 Liesen과 Strakoš의 책\cite{liesen2013krylov}이 있다.

\chapterreferences

